\section{Related Work}
\label{sec:related}

\paragraph{Tool-using agents and benchmarks.}
Tool-using LLM agents are evaluated on a small, fairly stable set of
suites. The Berkeley Function-Calling Leaderboard~\citep{bfcl2024,bfclv4_2025}
started as single-turn API selection and has grown into a stateful
multi-turn suite. $\tau$-bench~\citep{taubench2024} pushes in a different
direction, scoring dialog-driven retail and airline agents that must
reconcile a user goal against a tool catalog. Both report
\emph{end-to-end task success}, which fuses tool selection with reasoning,
parsing, and policy errors into one number. Trace-level harnesses
\citep{appworld2024,reactagent2023} expose more, yet still report
pass-rate as the headline. This is the wrong granularity for the
failure we care about; TFR isolates a portable slice of it across
vendors and benchmarks.

\paragraph{Empirical analysis of agent failures.}
Practitioners have informally reported that frontier models call tools
that were never declared. \citet{czapla2026} reproduces the behavior
across Claude, Gemini, and Grok, noting that ``Claude 4.5 hallucinated
access to a tool I hadn't given it yet \ldots I've reproduced similar
behavior with Gemini and Grok.'' On the academic side,
\citet{bugstudy2026} mine 998 bug reports from CrewAI and LangChain
and find ``API misuse,'' ``API incompatibility,'' and ``Documentation
Desync'' dominating the root-cause distribution, with the Self-Action
lifecycle stage hit hardest. Neither line puts a number on the failure
under controlled conditions. We do, formalizing the unauthorized-tool
subclass as TFR and reporting it across five models.

\paragraph{Prior tool-hallucination metrics.}
Tool-call name errors have been measured under several aggregations,
and the choice of denominator matters. \citet{li2023apibank} introduce
``API Hallucination'' in API-Bank as a ground-truth-vs-prediction
name-mismatch error category, with per-task error rates up to $61.4\%$
on then-current models. The closest prior measurement to ours is
\citet{huang2024metatool}, whose MetaTool sub-task~3 deliberately
removes the correct tool from the registry and tracks the resulting
fabrication through a per-query Correct Selection Rate. That isolates the
registry-membership phenomenon, but stays at the per-query
level on single-turn instructions. \citet{cao2025reliability} name the
phenomenon directly (``tool type hallucination,'' where the model
fabricates a tool absent from the available set), but their fix is a
training-time \emph{reliability alignment} that expands the action space
(defer / clarify / switch), and they report it inside a composite
reliability score rather than as a standalone per-call rate. We share
their target failure but differ on both axes: a per-call rate, and a
training-free inference-time intervention. The ``does the called tool
exist'' question is not itself new: ToolBeHonest~\citep{toolbehonest2024}
catalogues a \emph{non-existent tool} error category qualitatively, and
BFCL v4's hallucination metric \citep{bfclv4_2025} measures a different
failure, invoking a tool when none is appropriate. What is new is the
per-call operationalization across capability tiers and vendors, and the
content-controlled ablation (Section~\ref{sec:mechanism}) it makes
possible: TFR extends MetaTool's premise from \emph{per-query} to
\emph{per-call} on multi-turn agentic traces, and from controlled
removal to the natural distribution of registry-membership violations.
Per-call is what lets the intervention in Section~\ref{sec:method:rvr}
be evaluated against a content-controlled $C_{0.7}$ baseline, a contrast
aggregate metrics cannot expose. Appendix Table~\ref{tab:priorart}
positions TFR against these prior measurements by what each one counts;
the denominators and elicitation setups differ, so the numbers are not
directly comparable, but the granularity gap is the point.

\paragraph{Constrained decoding and self-correction.}
A separate line constrains agents at decode time. Grammar-constrained
and JSON-mode decoding \citep{outlines2023,jsonmode2024} prevent
malformed calls. Provider-side function-calling APIs reject unknown
names but surface the error as an opaque exception
\citep{openaitools2024,anthropictools2024}. Self-correction methods
such as Reflexion \citep{reflexion2023} and self-debug
\citep{selfdebug2023} push the diagnosis back onto the agent's own
trace. The closest neighbor in method shape is CRITIC
\citep{gou2024critic}, an external-feedback verify-then-correct loop
that needs no training and treats the model as a black box. RVR is a
registry-membership-specialized instance of that pattern. What
distinguishes it is the content-controlled $C_{0.5}/C_{0.7}/C_1$
ablation, which separates whether the gain comes from retrying, from
a structured-error code, or from the registry list itself. Production
frameworks have started to adopt structured tool-not-found feedback;
the LangChain community \citep{langchain_toolnotfound2025} explicitly
proposes returning a structured \texttt{ToolMessage} on a failed call,
but stops short of echoing the available registry. That gap is exactly
what our $C_{0.7}/C_1$ contrast measures.

\paragraph{Closely related interventions.}
Three concurrent works deserve direct comparison. \citet{licl2026}
introduce Localized In-Context Learning (L-ICL) for symbolic planners,
locating the first constraint violation and injecting a minimal
corrective example, lifting valid-plan rate from $59\%$ to $89\%$ on
$8{\times}8$ gridworld. RVR is similarly minimal but operates on
tool-using agents and pulls its signal from the runtime registry rather
than curated examples. \citet{fissiongrpo2026} train Qwen3-8B with
Fission-GRPO, an RL recipe that resamples failed trajectories with
diagnostic feedback, gaining $+5.7$pp error recovery on BFCL-v4
multi-turn. That is a training-time fix on the same benchmark family
we target at inference time, and the two are orthogonal.
\citet{englander2026} describe a related failure: agents discover
useful signals but fail to act on them (``agents use the environment
to fetch expected information, but not to revise their strategy'').
RVR addresses the dual problem of acting on \emph{unobserved} tools,
evaluated as a paired content-controlled contrast rather than
observationally.

\paragraph{Adjacent work.}
\citet{stableTOOLbench2024} targets tool-server stability rather than
name existence, and production tool-skill packaging
\citep{anthropicagentskills2025} is orthogonal.

\paragraph{Concurrent work on tool hallucination (2026).}
Three 2026 papers bracket our setting. \citet{yin2026reasoningtrap}
introduce SimpleToolHalluBench, whose ``no tool available'' and
``only distractor tools'' regimes mirror our target-removed probe, and
report that strengthening reasoning (RL/SFT/prompting) \emph{monotonically}
amplifies tool hallucination. \citet{qi2026brief} find a \emph{non-monotonic}
effect on a different axis, chain-of-thought \emph{budget} at fixed scale
on Qwen2.5-1.5B function calling. Our axis is orthogonal to both: model
\emph{scale} at fixed (disabled) thinking, where the non-monotonicity
appears across the Qwen3 size ladder rather than along a reasoning knob.
\citet{healy2026internal} detect tool-selection hallucination from a model's
internal representations in a single forward pass; we work purely
behaviorally and, rather than detect, \emph{prevent} with a message-level
re-prompt that needs no logit or activation access.
